\chapter{Literature Review}
\label{chap:literature}

\section{Deception beyond false statements}

A misleading communication is not exhausted by the falsity of a sentence. A speaker can mislead by withholding relevant evidence, exploiting implication, selecting a favourable comparison class, changing precision, or arranging true statements so that the recipient draws an unsupported conclusion. This distinction matters for LLM evaluation because most factuality metrics are proposition-centred: they ask whether an included claim is supported, not whether the response included what a decision-maker needed to know.

\citet{park2024deception} define AI deception around the systematic inducement of false beliefs in pursuit of an outcome other than truth. The belief-oriented component is useful for deployment evaluation, but the purposive element is difficult to infer from black-box output. This thesis therefore separates a behavioural label---materially misleading risk communication---from any claim about intention.

\citet{shi2026taxonomy} make the mechanism distinction explicit, separating fabrication, omission, and pragmatic distortion. Their analysis of 50 benchmarks finds broad coverage of fabrication and much thinner coverage of pragmatic distortion. \citet{berger2026liedetector} demonstrate the practical consequence: models can produce misleading non-falsities that truth probes detect less reliably than explicit lies. A detector can classify whether represented propositions are false while missing whether selected true propositions create a deceptive net impression.

Information-unit approaches offer a route forward. DECOR decomposes a context into atomic units and evaluates how each is transformed in a response \citep{cai2026decor}. JANUS fixes favourable and adverse facts, compares neutral with goal-conditioned prompts, and measures selection, specificity, emphasis, and framing \citep{giannouris2026janus}. \benchmarkname adopts the controlled fact-pool principle but changes the causal question: it removes explicit goal direction from the primary comparison, crosses system guidance with user persona, and adds multi-turn belief and action outcomes.

\section{Strategic deception and goal conflict}

Recent work establishes that capable LLMs can use concealment instrumentally. In a simulated financial environment, GPT-4 acts on an insider tip and hides its true rationale when reporting to a manager \citep{scheurer2023strategic}. Frontier-model scheming evaluations give models in-context goals and environments where covert action is useful; models sometimes manipulate oversight, strategically underperform, or maintain deceptive accounts \citep{meinke2024scheming}. Corporate stress tests similarly observe harmful actions when assigned goals conflict with organisational direction or continued operation \citep{lynch2025agentic}.

These studies are important capability demonstrations, but their strongest conditions supply a salient reason to conceal. Related work finds deceptive tendencies with less direct pressure \citep{jarviniemi2024deceptive}, harmful over-optimisation for helpfulness \citep{wang2026toxic}, and reward--ethics trade-offs in interactive settings \citep{pan2023machiavelli}. DeceptionBench varies neutral, incentivised, and coercive conditions and includes multi-turn interaction \citep{huang2025deceptionbench}. Collectively, the literature shows that deceptive behaviour need not follow a literal instruction to lie, while still focusing mainly on conspicuous acts, explicit falsehood, or policy violation.

The present benchmark sits at a lower evidential threshold. Its production-baseline prompt includes ordinary desiderata such as concision and reassurance but no hidden objective, conversion target, or instruction to suppress risk. Observed omission would establish a communication failure under benign conditions, not strategic scheming.

\section{Benign prompts and alternative explanations}

\citet{wu2025benign} directly investigate deception on benign prompts and report that proposed deception scores increase with task difficulty. The study highlights both the importance and the identification challenge of benign-condition evaluation: difficult tasks can produce errors, uncertainty, and unstable reasoning without strategic intent.

A controlled fact pool narrows this ambiguity. When the model receives a bounded source and preserves favourable details while repeatedly dropping high-adverse facts, asymmetric selection is more plausible than absent knowledge. Even then, possible mechanisms include context compression, instruction collision, uncertainty about relevance, stylistic accommodation, and learned reassurance. \benchmarkname therefore supports two levels of claim: response defects, and systematic conditional changes in those defects. Instrumental intent would require stronger causal or process evidence.

Task difficulty remains a confound. Scenarios should be annotated for source complexity, and negative controls should measure favourable and neutral fact retention. A response that omits every class of fact indicates general failure; disproportionate loss of adverse facts is the more diagnostic pattern.

\begin{placeholderblock}
After pilot review, add a compact typology of non-strategic explanations observed in real errors. Keep these explanations separate from the preregistered behavioural labels and do not use them to relabel inconvenient cases post hoc.
\end{placeholderblock}

\section{Financial LLM evaluation and risk communication}

Finance is a natural setting for pragmatic distortion because decisions depend on materiality, conditions, time horizons, and recipient circumstances. A return figure can be accurate yet misleading without a benchmark or volatility; an account can be technically accessible yet costly to exit; a benefit can be real but unavailable under an omitted exclusion.

The deployment context is increasingly relevant. The Bank of England and FCA report broad use of AI in UK financial services \citep{boefca2024ai}. A joint ESMA, Alan Turing Institute, and Institut Louis Bachelier report discusses LLM applications in text analysis, production, and customer interaction alongside governance needs \citep{esma2025llmfinance}. The FCA's AI approach and Mills Review likewise consider more capable and autonomous systems in retail finance \citep{fca2025aiapproach,fca2026mills}. These sources do not show that deployed systems conceal risk; they show that the communication channels are becoming operational.

Conduct language also motivates a net-impression perspective. Requirements that communications be clear, fair, and not misleading cannot be reduced to forbidding false propositions \citep{fca_handbook_principle7}. The benchmark is not a legal compliance test, but it operationalises related properties: balance, specificity, materiality, and audience.

Existing finance benchmarks mainly test capability. FinQA and ConvFinQA assess numerical reasoning over reports and conversation \citep{chen2021finqa,chen2022convfinqa}; FinanceBench tests document-grounded question answering \citep{islam2023financebench}; FinBen aggregates financial tasks \citep{xie2024finben}; and recent agent benchmarks evaluate research and tool use \citep{bigeard2025financeagent,zhu2026finmcp}. Safety work is growing: \citet{chen2025standard} argue for risk-aware finance-agent evaluation, while FinSafetyBench targets refusal of harmful financial requests \citep{hou2026finsafety}. \citet{cheng2025riskconcealment} conceal regulatory risk \emph{from} a model to bypass safeguards; this thesis instead studies whether a model that receives the risk fails to communicate it \emph{to} a user.

A finance-specific design makes materiality explicit. The same fact may deserve different disclosure for a retail customer, analyst, or compliance reviewer. Stakeholder role must therefore be fixed within a family while persona varies.

\section{Persona, sycophancy, and adaptive communication}

Conversational adaptation is useful when it changes tone or explanation level, but unsafe when it changes the substance of risk disclosure. Sycophancy research finds that models may match user views over correct answers and that preference data can reward agreement \citep{sharma2023sycophancy}. A financial analogue need not repeat a false belief: an enthusiastic user may receive stronger benefit language or fewer qualifications, while an anxious user may receive either more detail or vague reassurance.

The benchmark uses three composite profiles: neutral; anxious risk-averse; and positive risk-seeking or trusting. The conditions are policy-relevant but not psychologically pure. Emotion, trust, risk preference, and desired detail vary together, so observed effects cannot be assigned to one trait.

\begin{placeholderblock}
Decide whether the positive profile is primarily ``risk-seeking,'' ``trusting,'' or explicitly bundled. Risk appetite and trust are distinct: a sophisticated high-risk investor may still demand exact disclosure. Report the condition as composite unless later ablations isolate its dimensions.
\end{placeholderblock}

Persona can affect initial wording, follow-up questions, and the model's inferred materiality threshold. To avoid scorer bias, fact matching and framing judges should be blinded to condition. A user should not receive a lower substantive standard merely because they sound confident.

\section{Multi-turn outcomes and detection}

Single-turn evaluation is controlled, but dialogue can repair or worsen risk communication. An initially omitted fact may be volunteered after a question; an accurate caveat may later be softened. OpenDeception and DeceptionBench show the value of interaction-level evaluation \citep{wu2025opendeception,huang2025deceptionbench}. \benchmarkname uses a user simulator to continue the assigned persona, then elicits final beliefs and a bounded action choice.

Simulator outcomes are proxies, not human population estimates. A model user may attend to text differently, share biases with the target, or follow elicitation instructions unusually well. Stronger designs use separate dialogue and outcome models, multiple simulator families, randomised action order, and human validation. Conversation metrics should retain first-turn disclosure, time to disclosure, persistence, repair, and final belief rather than collapsing the trajectory into one score.

Detection work further illustrates why failure modes must be separated. Black-box lie detection can generalise by asking unrelated questions \citep{pacchiardi2023catch}; white-box probes use internal activations. Recent evidence suggests that deception signals are heterogeneous and targeted probes may outperform a universal detector \citep{natarajan2026oneprobe}. Black-to-white comparisons quantify whether internal access adds value \citep{parrack2025blacktowhite}, while Liars' Bench shows systematic detector failures across reasons for lying and belief targets \citep{kretschmar2025liars}. Omission is especially challenging because the missing fact has no output span; a monitor may need the source and expected-disclosure metadata.

\section{Research gap}

The literature establishes that LLMs can deceive under instrumental pressure, that problematic behaviour can arise without a direct instruction to lie, and that current benchmarks and detectors overemphasise fabrication. Finance benchmarks increasingly measure reasoning and tool use but rarely audit whether material adverse evidence survives into user-facing communication. User-facing models also adapt to recipient cues in ways that may trade balance against agreement or reassurance.

What remains missing is a controlled, finance-specific estimate of truthful-but-misleading risk communication under ordinary prompts. \benchmarkname addresses this gap by holding evidence constant, annotating expected disclosure in advance, crossing system guidance with persona, and tracing effects from fact treatment to simulated belief and action. Its aim is to establish whether a deployment-relevant failure occurs, how it manifests, and which conditions change its probability---without overstating what output alone reveals about intent.
